\documentclass[12pt, letterpaper]{article} 

\usepackage{amsmath, amssymb, amsthm, enumerate}

\title{MATH 112A - Mathematical Biology}
\author{Dani Nguyen}
\date{Winter 2025}

\begin{document}
\maketitle
\tableofcontents

\section{Review}
\subsection{Ordinary Differential Equations}
An ODE can be written as $F(y,y',y'', ..., y^{(k)})=0$. The order of the ODE is the order of the highest derivative. \\

\noindent Linear ODEs are a special class of ODEs that have the form:
\begin{align*}
    a_0(t)y+a_1(t)y'+...+a_k(t)y^{(k)} = f(t)
\end{align*}

\noindent There exists two special cases:
\begin{enumerate}
    \item If $f(t) = 0$ then the linear ODE is called homogeneous.
    \item If $a_i(t)$ are constant and $f(t)=0$ then we have a linear homogeneous ODE with constant coefficients.
\end{enumerate}
\boxed{ex} \\
$y' = 0$ \\
This is a first order linear homogeneous ODE with constant coefficients. \\ \\
$y'''+ \sin t\cdot y'' + (t^2+1)y+20t =0$ \\ 
This is a third order linear non-homogeneous ODE \\ \\ 
$y'\cdot y +20$ \\ 
This is a first order non-linear non-homogeneous ODE. \\ \\
$y^2+y''=0$ \\
This is a second order non-linear homogeneous ODE. \\ \\

\noindent The simplest type of ODEs is linear, homogeneous, with constant coefficients.\\
\boxed{ex}\quad $y''+y' -6y = 0$  \\
Look for solutions as:
\begin{align*}
    y &= e^{pt}\\
    y' &= pe^{pt}\\
    y'' &= p^2e^{pt}
\end{align*}
\begin{align*}
    p^2e^{pt} + pe^{pt} - 6e^{pt} &=0 \\
    p^2 + p -6 &= 0 \\
    p &= \frac{-1 \pm 5}{2} \\
    &= -3 \text{ or } 2
\end{align*}
The general solution is $y(t) = Ae^{-3t} + Be^{2t}$
Initial value problem:
Suppose
\begin{align*}
    y(0) &= 1 \\
    y'(0) &= 0 
\end{align*}
Then use the general solution to find $A$ and $B$ \\
$y' = -3Ae^{-3t} + 2Be^{2t}$ \\
$\begin{cases}
    y(0) = Ae^{-3\cdot 0} + Be^{2\cdot 0} = A+B = 1 \\ 
    y'(0) = -3A+ 2B = 0
\end{cases}$ \\ 
Solving the linear system gives $A = \frac{2}{5}$ and $B = \frac{3}{5}$ \\ 

\noindent \boxed{ex}\quad $y'+2y-e^t = 0 $ 

$p(t) = 2$ \quad $q(t) = e^t$ \\
$E(t) = \exp \int_{0}^{t}p(t)dt = e^{2t}$ \\ 
$e^{2t}y'+2e^{2t}y =  e^{2t}\cdot e^{t}$ \\
$\frac{d}{dx} e^{2t}y = e^{3t}$ \\
$e^{2t}y = \frac 13 e^{3t} + c$ \\
$y = \frac{1}{3} e^t +ce^{-2t}$ \\ \\ 

\noindent Separable ODEs
\begin{align*}
    \frac{dy}{dt} &= g(y)h(t) \\
    \frac{dy}{g(y)} &= h(t)dt \\
    \int \frac{dy}{g(y)} &= \int h(t)dt 
\end{align*}

\noindent A more realistic model of population growth where $r$ = growth rate and $k$ = carrying capacity.
\begin{align*}
    \frac{dy}{dt} = ry(1-\frac yk)
\end{align*}
When $y<<k$, RHS $\approx ry$

\end{document}